\documentclass[12pt]{article}
\usepackage[utf8]{inputenc}
\usepackage[T1]{fontenc}
\usepackage{hyperref}
\title{Tabasaran Language Resources for K-12 Curriculum Development}
\author{Compiled by Claude Sonnet 5 (claude-sonnet-5) for Andr\'as Kornai}
\date{}
\begin{document}
\maketitle

\section{Introduction}
Tabasaran (ISO 639-3 \texttt{tab}, Glottocode taba1259) is a Nakh-Daghestanian language of Dagestan, Russia, with about 130,000 speakers. It has been written in Cyrillic since 1938. The search found dictionaries and a grammar sketch. It found no school textbooks, and a search in Russian would be the next step.

\section{Primary Educational Resources}
No school textbooks or curriculum materials turned up. The Webonary dictionary is described as a tool for speakers learning to read and write Tabasaran.

\section{Linguistic and Reference Materials}
\begin{itemize}
\item Tabasaran--Russian dictionary (compiled by Sabrina Shikhalieva, 2019) on SIL Webonary: \url{https://www.webonary.org/tabasaran/overview/introduction/}
\item Freelang Tabassaran--English word list (about 985 entries, romanized): \url{https://www.freelang.net/dictionary/tabassaran.php}
\end{itemize}

\section{Religious and Community Resources}
The search found none.

\section{Digital and Audio Resources}
The Webonary and Freelang dictionaries above are online. The search found no audio.

\section{Curriculum Development Context}
The language has a literary standard and a Cyrillic orthography. The first steps would be to locate existing Dagestani school textbooks, which this search did not find, and to build on the Webonary dictionary.
\end{document}